\begin{answer}
The implementation in \texttt{src/k\_means/k\_means.py} proceeds as follows:

\textbf{Initialization} (\texttt{init\_centroids}): For each of the $k=16$ cluster centroids, we randomly select a pixel from the image (uniformly at random) and use its $(r,g,b)$ values as the initial centroid.

\textbf{Centroid update} (\texttt{update\_centroids}): We run the standard K-means update loop for at most \texttt{max\_iter} iterations. At each iteration we (1) assign every pixel to its nearest centroid (minimizing squared Euclidean distance in RGB space), then (2) recompute each centroid as the mean of all pixels assigned to it. We stop early if the centroids do not change between iterations.

\textbf{Image update} (\texttt{update\_image}): We replace every pixel in the large image with the value of the closest centroid (learned on the small image).

The compressed image and original image are shown below (generated by running the script):

\begin{center}
\textit{[Original peppers-large.tiff and compressed version with $k=16$ colors shown side by side.]}
\end{center}

The compressed image visually resembles the original, with visible colour quantization artefacts on smooth gradient regions, confirming the implementation is correct.
\end{answer}
